\chapter{Projekt rozwiązania i~implementacja}
\label{chap:sol_and_impl}
Projekt rozwiązania jest ważnym elementem do rozważenia,
by implementacja nie zawierała błędów.
Poniżej zostaną przedstawione projekt rozwiązania,
a~następnie implementacja.

\section{Struktura projektu i~przebieg optymalizacji}
\label{sec:proj_struct}
Jak oceniono w~\cref{chap:prob_and_req},
biblioteka \gls{GAAML} powinna
udostępniać dwie funkcje.
Aby wykonać to w~sposób umożliwiający
łatwe modyfikowanie i~rozwijanie
projekt został podzielony na części.
Najważniejszą jest główna funkcja
przeprowadzająca pełną ewolucję,
używa ona wszystkich niżej wymienianych
elementów bezpośrednio lub pośrednio.

Pierwszą czynnością, którą wykonuje
funkcja po uruchomieniu,
jest przygotowanie danych do uczenia
na nich i~sprawdzania dopasowania.
Następnie tworzy obiekt pokoleń,
który ma za zadanie zarządzać całym
biegiem ewolucji uruchamiając
ją na kolejnych populacjach
oraz zapisywanie ich statystyki.
Po stworzeniu go zostaje oddelegowane
mu zadanie przeprowadzenia ewolucji,
a~po zakończeniu tego procesu główna
funkcja odbiera wyniki i~wyświetla je użytkownikowi,
jeżeli na to wskazuje flaga rysowania wykresów.

Przygotowanie polega wpierw na sprawdzeniu
czy ilość atrybutów jest równa między danymi treningowymi,
testowymi oraz walidacyjnymi, jeżeli są obecne.
Następnym jest sprawdzenie danych czy dotyczą
problematyki klasyfikacji, czy regresji.
Oraz czy jest to zgodne z~deklaracją
przy wywołaniu funkcji.
Ostatecznie definiowana jest funkcja przyjmująca
tylko osobnika i~ścieżkę do zapisu o~nim informacji.
Funkcja ta tworzy model, uczy go i~oblicza
błąd dopasowania jego odpowiedzi do tych
przewidywanych w~danych testowych oraz walidacyjnych,
a~jeżeli te drugie są niedostępne to treningowych.
\gls{MSE} będzie używany do liczenia błędów dopasowania.
Drugi z~policzonych błędów będzie zwany
dalej błędem dopasowania walidacyjnego,
a~pierwsze analogicznie błędem dopasowania testowego.
Błąd oraz uzyskane po korekcji przystosowanie,
z~danych testowych nie będą wykorzystywane
do żadnego procesu ewolucji,
pozwalają jedynie na ujrzenie pełnego obrazu,
jak się ta wartość przystosowania zmieniała w~czasie.
Następnie przed zwróceniem obydwu dopasowań
zapisuje wagi nauczonego modelu do pliku
by użytkownik mógł je potem w~łatwy
sposób użyć we własnym modelu,
a~do drugiego pliku wpisywane
są następujące metadane: liczba epok,
wielkość wsadu, backend, liczba warstw ukrytych,
historia strat podczas kolejnych epok nauki,
błąd dopasowania walidacyjnego
oraz błąd dopasowania testowego.

\section{Klasa pokoleń}
\label{sec:generations_class}
Obiekt pokoleń jest tworzony z~klasy,
która zapisuje statystykę i~pełny obraz
przystosowania z~populacji
oraz posuwa proces ewolucji do przodu
poprzez wywoływanie kolejnych generacji.
Powinien umożliwiać zobaczenie obecnego numeru pokolenia.
Powinien również umożliwiać zobaczenie zapisanej
statystyki oraz wszystkich obliczonych przystosowań.

Tworzy na początku swojego działania obiekt populacji.
Następnie po dostaniu polecenia przeprowadzenia
ewolucji sprawdza ile generacji już było.
Jeżeli jeszcze nie została osiągnięta zadana ilość
jest wywoływana metoda następnej generacji
z~obiektu populacji z~numerem,
obliczanego w~danym momencie, pokolenia.
Następnie pobiera wyniki z~populacji i~je zapisuje.

\section{Klasa populacji}
\label{sec:population_class}
Obiekt populacji jest tworzony z~klasy.
Powinien umożliwiać pobranie statystyk jaki
jest najlepszy osobnik i~najgorszy z~populacji
oraz ich wartości wraz ze średnią przystosowań.
Wartości te są obliczane przez wywoływanie
funkcji przygotowanej na początku całego procesu.
Jest ona uruchamiana tyle razy dla każdego osobnika
ile zostało ustalone przez liczbę obliczeń przystosowania.
Obiekt populacji ma podawać do tej
funkcji unikalne ścieżki do folderów,
które będą się znajdowały wewnątrz głównego
folderu podanego uprzednio przez użytkownika.

Podczas tworzenia obiektu od razu generowana
jest populacja początkowa i~obliczane są wartości
przystosowania osobników należących do niej.
Generowanie osobników do populacji
początkowej składa się z~dwóch etapów.
Pierwszym z~nich jest wygenerowanie części sekwencyjnej,
które następuje przez wylosowanie zadanej liczby bitów.
W~podobny sposób w~drugim etapie
generowane są części listowe,
jednakże liczba bitów do wylosowania jest
określana dzięki parametrowi liczby warstw
ukrytych z~uprzednio zdobytej części sekwencyjnej.
Obliczenie przystosowania polega na nauczeniu
kilku sieci o~tej samej architekturze,
obliczeniu dla nich błędu dopasowania walidacyjnego,
przekształceniu ich na przystosowanie walidacyjne
za pomocą funkcji korygującej,
a~ostatecznie uśrednieniu uzyskanych wyników.
Po wywołaniu metody następnej generacji
jest przeprowadzany jeden pełny obieg
wnętrza pętli pokazanej
na \cref{fig:genetic_flow},
czyli selekcja, krzyżowanie,
mutacja oraz obliczenie przystosowania.
Populacja powinna zawierać osobniki jednoznacznie
reprezentujące modele sieci neuronowych
i~udostępniające metody takie jak mutacja,
wygenerowanie punktu/punktów do krzyżowania,
krzyżowanie, zapisanie do pliku
oraz odczytanie z~pliku.

\section{Kod genetyczny osobnika}
\label{sec:individual_genetic_code}
W~niniejszym dokumencie kod genetyczny
będzie reprezentował architekturę sieci
i~jej parametry uczenia.
Aby zapewnić łatwość modyfikacji
i~rozwijania rozwiązania zostanie
on podzielony na trzy podosobniki.
Każdy z~nich będzie odpowiadał
za jedną z~części genotypu.
Wspomniane części mają być reprezentowane
przez kod binarny i~mają to być:
\begin{nospitemize}
  \item Część sekwencyjna:
    \begin{nospitemize}
      \item liczba warstw ukrytych:
        5 bitów $\mathbb{Z} \cap \left[1, 32\right]$,
      \item ziarno generatora liczby neuronów:
        10 bitów $\mathbb{Z} \cap \left[0, 1023\right]$,
      \item ziarno generatora typów aktywacji:
        10 bitów $\mathbb{Z} \cap \left[0, 1023\right]$,
      \item optymalizator:
        1 bit $\mathbb{Z} \cap \left[0, 1\right]$
        (%
          0 {\textrightarrow} SGD,
          1 {\textrightarrow} Adam%
        ),
      % 2^(-14)*[1, 2^(12)]==[2^(-14) - .25]
      \item wskaźnik uczenia się (learning rate):
        12 bitów $2^{-14}\left(
          \mathbb{Z} \cap \left[1, 2^{12}\right]
        \right)$,
      \item liczba epok:
        7 bitów $\mathbb{Z} \cap \left[1, 128\right]$,
      \item wielkość partii:
        10 bitów $\mathbb{Z} \cap \left[1, 1024\right]$,
    \end{nospitemize}
  \item Części listowe
    -- generowane po części sekwencyjnej;
    minimalne i~maksymalne ich długości
    są zależne od minimalnej
    i~maksymalnej wartości
    dla liczby warstw ukrytych:
    \begin{nospitemize}
      \item lista ilości neuronów w~warstwach ukrytych:
        \\\hspace*{1em}liczba elementów -- taka sama:
          $\mathbb{Z} \cap \left[1, 32\right]$,
        \\\hspace*{1em}zawiera liczby z~zakresu:
          $\mathbb{Z} \cap \left[2, 724\right]$,
        \\\hspace*{1em}każdy element
          składa się z~10 bitów,
      \item lista typów aktywacji
        warstw ukrytych i~wyjściowej:
          \\\hspace*{1em}liczba elementów -- o~$1$ większa
            (uwzględnia warstwę wyjściową):
            $\mathbb{Z} \cap \left[2, 33\right]$,
          \\\hspace*{1em}zawiera liczby z~zakresu:
            $\mathbb{Z} \cap \left[0, 3\right]$
            (%
              $0$ {\textrightarrow} relu,
              $1$ {\textrightarrow} tanh,
              $2$ {\textrightarrow} sigmoid,
              $3$ {\textrightarrow} leaky relu%
            ),
          \\\hspace*{1em}każdy element
            składa się z~2 bitów.
    \end{nospitemize}
\end{nospitemize}

Zostało przyjęte ograniczenie by macierz
wag połączeń w~warstwach ukrytych sieci
była niewiększa niż $128\text{MB}$.
Każda komórka w~macierzy ma typ \texttt{float}.
Typ ten może zajmować w~zależności
od implementacji zajmuje $64\text{b} = 8\text{B}$
lub $32\text{b} = 4\text{B}$.
Jeżeli przyjęta zostanie wersja z~typem
\texttt{float64} oraz maksymalna liczba
warstw będzie ustalona jako $32$, to:
\begin{equation}
  \begin{gathered}
    P_{\max}
      = 128\,\text{MB}
      = 131072\,\text{kB}
      = 134217728\,\text{B}\\
    % ---
    \Downarrow\\
    % ---
    n_{\max}^2
      \leqslant P_{\max}
        \times \frac{1}{8}\frac{\text{komórek}}{\text{B}}
        \times \frac{1}{32}\frac{1}{\text{warstwę}}
      = 524288\frac{\text{komórek}}{\text{warstwę}}\\
    % ---
    \Downarrow\\
    % ---
    \sqrt{524288}
      \approx 724.077343935
      \geqslant n_{\max} = 724
  \end{gathered}
\end{equation}
Oznacza to, że dla maksymalnej liczby
neuronów w~każdej macierzy połączeń
może być maksymalnie $524288$ komórek.
Czyli maksymalna liczba neuronów
w~warstwach musi być mniejsza niż
pierwiastek kwadratowy wspomnianej liczby,
więc maksymalna liczba neuronów w~warstwie to $724$.
Minimalna wartość ($2$) wynika z~faktu,
że architektura z~pojedynczym neuronem
w~dowolnej warstwie ukrytej działałaby
w~sposób zbliżony do modelu liniowego.

Ważnym mechanizmem podczas czytania kodu genetycznego
w~celu utworzenia fenotypu są ziarna generatora.
Jeżeli liczba warstw ukrytych z~części sekwencyjnej,
wynosiłoby więcej niż liczba elementów w~listach,
to~należy uzupełnić listy ilości neuronów
i~typów aktywacji do danej liczby elementów
do czego zostanie użyty generator z~odpowiednim ziarnem.
Jeżeli liczba warstw ukrytych jest
mniejsza niż liczba elementów w~listach,
to nadmiarowe elementy zostaną zignorowane w~fenotypie.

\begin{figure}
  \centering
  \resizebox{\textwidth}{!}{%
    \begin{tikzpicture}[
  font=\sffamily\small,
  umlclass/.style={
    draw=blue!55!black,
    thick,
    rounded corners=1.5pt,
    fill=blue!4,
    inner sep=0pt,
    align=left,
  },
  inheritance/.style={
    draw=blue!55!black,
    thick,
    -{Triangle[open,length=3.4mm,width=4.2mm]},
  },
  association/.style={
    draw=gray!70!black,
    thick,
    -{Stealth[length=2.7mm,width=2.1mm]},
  },
]
  % UML class nodes
  \node[umlclass] (Ind) at (0,0) {%
    \begin{tabular}{>{\raggedright\arraybackslash}p{12cm}}
      \multicolumn{1}{c}{%
        {\textless\textless}%
        abstract%
        {\textgreater\textgreater}%
      }
      \\\multicolumn{1}{c}{
        \textbf{Individual{\textless}CPT, T{\textgreater}}
      }
      \\\multicolumn{1}{c}{as \textbf{(I)}}
      \\\hline\hspace{0.1em}+ T genotype
      \\\hline\hspace{0.1em}\textit{+ mutate(): None}
      \\\hspace{0.1em}\textit{+ get\_crosspoint(%
        Individual a,
        Individual b%
      ): CPT}
      \\\hspace{0.1em}\textit{+ crossover(%
        Individual a,
        Individual b,
        CPT crosspoint%
      ): tuple{\textless}Individual, Individual{\textgreater}}
      % \\\hline \hspace{0.1em}+ int ala
    \end{tabular}%
  };

  \node[
    umlclass,
    below left=1.5cm and -3.1cm of Ind,
  ] (SI) {%
    \begin{tabular}{>{\raggedright\arraybackslash}p{4.1cm}}
      \multicolumn{1}{c}{\textbf{Sequential Individual}}
      \\\multicolumn{1}{c}{as \textbf{(SI)}}
      \\\hline \hspace{0.1em}
      \\\hline \hspace{0.1em}
    \end{tabular}%
  };

  \node[
    umlclass,
    below right=1.5cm and -3.1cm of Ind,
  ] (LI) {%
    \begin{tabular}{>{\raggedright\arraybackslash}p{3.5cm}}
      \multicolumn{1}{c}{\textbf{List Individual}}
      \\\multicolumn{1}{c}{as \textbf{(LI)}}
      \\\hline \hspace{0.1em}
      \\\hline \hspace{0.1em}
    \end{tabular}%
  };

  \node[
    umlclass,
    below=1.0cm of SI,
  ] (MNSI) {%
    \begin{tabular}{>{\raggedright\arraybackslash}p{6.25cm}}
      \multicolumn{1}{c}{\textbf{Sequential Individual With Max Numbers}}
      \\\multicolumn{1}{c}{as \textbf{(MNSI)}}
      \\\hline \hspace{0.1em}
      \\\hline \hspace{0.1em}
    \end{tabular}%
  };

  \node[
    umlclass,
    below=1.0cm of LI,
  ] (MNLI) {%
    \begin{tabular}{>{\raggedright\arraybackslash}p{5.7cm}}
      \multicolumn{1}{c}{\textbf{List Individual With Max Numbers}}
      \\\multicolumn{1}{c}{as \textbf{(MNLI)}}
      \\\hline \hspace{0.1em}
      \\\hline \hspace{0.1em}
    \end{tabular}%
  };

  \node[
    umlclass,
    below=7cm of Ind,
  ] (NI) {%
    \begin{tabular}{>{\raggedright\arraybackslash}p{4.4cm}}
      \multicolumn{1}{c}{\textbf{Neural Network Individual}}
      \\\multicolumn{1}{c}{as \textbf{(NNI)}}
      \\\hline \hspace{0.1em}
      \\\hline \hspace{0.1em}
    \end{tabular}%
  };

  % UML inheritance
  \draw[inheritance]
    (NI.north)
    -- node[
      label,
      fill=white,
      align=center,
      yshift=-5.2mm,
      xshift=.4pt,
      font=\scriptsize,
    ] {\texttt{%
      {\textless}%
        tuple{\textless}int, lcpt, lcpt{\textgreater},%
    }\\\texttt{%
        tuple{\textless}MNSI, MNLI, MNLI{\textgreater}%
      {\textgreater}%
    }}
    (Ind.south)
  ;
  \coordinate (SI_inhP) at ($ (SI.north east|-Ind.south)+(3pt, 0) $);
  \draw[inheritance]
    (SI.north east)
    -- node[
      label,
      fill=white,
      yshift=-1.7mm,
      font=\scriptsize,
    ] {\texttt{{\textless}int, bytearray{\textgreater}}}
    (SI_inhP)
    % ($ (Ind.south)!.5!(Ind.south west) $)
  ;
  \draw[inheritance]
    (MNSI.north east)
    -- node[
      label,
      fill=white,
      yshift=-1.7mm,
      font=\scriptsize,
    ] {\texttt{{\textless}int, SI{\textgreater}}}
    ($ (Ind.south)!.5!(SI_inhP) $)
  ;
  \coordinate (LI_inhP) at ($ (LI.north west|-Ind.south)+(-3pt, 0) $);
  \draw[inheritance]
    (LI.north west)
    -- node[
      label,
      fill=white,
      yshift=-1.7mm,
      font=\scriptsize,
    ] {\texttt{{\textless}lcpt, bytearray{\textgreater}}}
    (LI_inhP)
    % ($ (Ind.south)!.5!(Ind.south east) $)
  ;
  \draw[inheritance]
    (MNLI.north west)
    -- node[
      label,
      fill=white,
      yshift=-1.7mm,
      font=\scriptsize,
    ] {\texttt{{\textless}lcpt, LI{\textgreater}}}
    ($ (Ind.south)!.5!(LI_inhP) $)
  ;

  % UML association
  \draw[association]
    (MNSI)
    --
    (SI)
  ;
  \draw[association]
    (MNLI)
    --
    (LI)
  ;
  \draw[association]
    (NI.west)
    to[out=170,in=-75]
    (MNSI.south)
  ;
  \draw[association]
    ($ (NI.north east)!0.667!(NI.south east) $)
    to[out=10,in=-105]
    ($ (MNLI.south west)!0.667!(MNLI.south east) $)
  ;
  \draw[association]
    ($ (NI.north east)!0.333!(NI.south east) $)
    to[out=20,in=-100]
    ($ (MNLI.south west)!0.333!(MNLI.south east) $)
  ;
\end{tikzpicture}%
  }
  {%
    \renewcommand{\thesubfigure}{I}%
    \phantomsubcaption\label{subfig:gen_parts_I}%
    \renewcommand{\thesubfigure}{SI}%
    \phantomsubcaption\label{subfig:gen_parts_SI}%
    \renewcommand{\thesubfigure}{LI}%
    \phantomsubcaption\label{subfig:gen_parts_LI}%
    \renewcommand{\thesubfigure}{MNSI}%
    \phantomsubcaption\label{subfig:gen_parts_MNSI}%
    \renewcommand{\thesubfigure}{MNLI}%
    \phantomsubcaption\label{subfig:gen_parts_MNLI}%
    \renewcommand{\thesubfigure}{NNI}%
    \phantomsubcaption\label{subfig:gen_parts_NNI}%
  }%
  \caption{
    Dziedziczenie i~powiązanie klas związanych
    z~osobnikiem representującym sieć neuronową (%
    \textbf{Neural Network Individual}
    \textbf{\subref*{subfig:gen_parts_NNI}}%
    ), dalej zwany osobnikiem.
    Dziedziczy on z~klasy abstrakcyjnej
    (ang. \textit{abstract}) osobnika (\textbf{Individual}
    \textbf{\subref*{subfig:gen_parts_I}}),
    jak i~inne osobniki, dalej zwane podosobnikami.
    Osobnik posiada jako genotyp
    zestaw trzech podosobników.
    Pierwszy jest typu podosobnika
    sekwencyjnego z~maksymalnymi liczbami
    (\textbf{Sequential Individual With Max Numbers}
    \textbf{\subref*{subfig:gen_parts_MNSI}}),
    reprezentuje on dane z~części sekwencyjnej.
    Dwa kolejne kodują dwie części listowe
    -- ilości neuronów i~typów aktywacji
    w~poszczególnych warstwach
    -- poprzez podosobnika list
    z~maksymalnymi liczbami
    (\textbf{List Individual With Max Numbers}
    \textbf{\subref*{subfig:gen_parts_MNLI}}).
    Podosobniki te mają za zadanie budowanie podfenotypu.
    Zarządzanie binarnym podgenotypem
    jest pozostawiane odpowiednio:
    podosobnikowi sekwencyjnemu
    (\textbf{Sequential Individual}
    \textbf{\subref*{subfig:gen_parts_SI}})
    i~listowemu (\textbf{List Individual}
    \textbf{\subref*{subfig:gen_parts_LI}}).
    One bezpośrednio implementują
    operatory krzyżowania i~mutacji.
    Typy punktów krzyżowania
    oznaczone jako \texttt{lcpt} to
    \texttt{tuple{\textless}int, int{\textgreater}}.
    Osobnik łączy dane z~podosobników
    i~aktualizuje fenotypy listowe.
  }
  \label{fig:gen_parts}
\end{figure}

Kod genetyczny ma strukturę zestawu trzech podosobników,
którzy zajmują się poprawnym krzyżowaniem poszczególnych
części jak pokazano na \cref{fig:gen_parts}.
Dzięki takiemu podziałowi kod powinien być łatwy
do zarządzania i~późniejszego rozwijania.

\section{Selekcja}
\label{sec:selection}
W~ogólnej sytuacji selekcja będzie odbywała się
metodą ruletki proporcjonalnej do przystosowania.
Selekcja musi jednakże odpowiadać na dwa szczególne
przypadki związane z~występowaniem zer
jako wartości przystosowań mianowicie:
częściowe występowanie i~całkowite.
W~drugim przypadku zostanie użyta selekcja losowa.
Natomiast w~pierwszym nadal będzie
używana selekcja ruletkowa,
jednakże z~uprzednim skorygowaniem wartości
przystosowań o~$a$~równe
$\displaystyle \frac{1}{100}\min_{i: v_i>0}{v_i}$,
gdzie $v$ to wektor wartości przystosowań osobników.
Jest to na przykład konieczne w~przypadku,
gdy wektor ma postać $\left[0,0,0,3.2157,0,0\right]$,
ponieważ do krzyżowania zostałby wzięty za każdym
razem jedyny osobnik z~niezerowym przystosowaniem.
Rozwiązanie to jest również korzystne gdy tylko
jeden z~osobników ma przystosowanie równe zeru,
ponieważ nie jest całkowicie
wykluczany z~procesu ewolucji.
Korekcja przystosowań powinna się odbywać
na wszystkich wartościach nie tylko tych zerowych,
ponieważ to by spowodowało zbyt duże zajęcie
przestrzeni ruletki przez osobniki,
które poprzednio w~ogóle jej niezajmowały.
Selekcja ruletkowa i~powyżej opisana metoda korekcji
została wybrana ze względu na możliwość
zachowania równowagi pomiędzy rozwijaniem
najlepszych znalezionych architektur,
a~przeszukiwaniem przestrzeni
w~poszukiwaniu nowych rozwiązań.
W~ten sposób osobniki o~wyższym przystosowaniu
są często wybierane do krzyżowania,
jednak również słabsze jednostki
mają szanse zostać wybrane.

\subsection*{Elitaryzm}
W~niniejszej pracy zrezygnowano z~elitaryzmu
w~celu zwiększenia różnorodności populacji
i~ograniczenia ryzyka przedwczesnej zbieżności,
a~najlepsze osobniki są zapisywane w~obiekcie
pokoleń opisanym w~\cref{sec:generations_class}.  

\section{Krzyżowanie}
\label{sec:crossover}
Krzyżowanie będzie odbywało się różnie dla części
sekwencyjnej i~listowej kodu genetycznego.
Choć obie składają się z~informacji binarnej,
to pierwsza z~nich ma stałą długość,
a~druga jest listą i~muszą być wzięte
pod uwagę różne długości podgenotypów oraz elementy list,
by nie było ``półelementów''.
Jednym ze wspólnych elementów w~tym procesie
jest fakt wykonywania operacji na binarnych podgenotypach,
które mają pozycje takie jak
ukazano na \cref{fig:bin_gen}.
W~każdym z~procesów uczestniczy dwóch
rodziców i~powstaje dwóch potomków.
Ostatnim elementem wspólnym jest podzielenie
operacji krzyżowania na dwie osobne podoperacje
losowania punktu/punktów przecięcia
oraz faktycznego wykonania cięcia
i~łączenia wedle podanych punktów.

\begin{figure}
  \centering
  \begin{tikzpicture}
  \def\orygbin{01011}
  \ExplSyntaxOn
    \def\foreachbit#1/#2 in #3#4{%
      \int_zero:N \l_tmpa_int
      \exp_args:Ne \tl_map_inline:nn {#3}
        {
          \edef#1{\int_use:N \l_tmpa_int}
          \edef#2{##1}
          #4
          \int_incr:N \l_tmpa_int
        }
    }
    \NewDocumentCommand{\strlen}{m}
    {
      \exp_args:Ne \tl_count:n {#1}
    }
  \ExplSyntaxOff

  \foreachbit \i/\bit in \orygbin{
    \node (b\i) at (\i em,0) {$\bit$};
  }
  \node[left=.3cm of b0] (b) {kod binarny:};

  \coordinate
    (h)
    at ($ (b0.south west)-(b0.north west) $)
  ;
  \coordinate
    (b1u)
    at ($ (b0.north east)!.5!(b1.north west) $)
  ;
  \coordinate
    (b2u)
    at ($ (b1.north east)!.5!(b2.north west) $)
  ;
  \coordinate
    (b0u)
    at ($ (b1u)!-1!(b2u) $)
  ;
  \foreachbit \i/\bit in {\orygbin}{{
    \ifnum\i<2
    \else
      \edef\lll{\number\numexpr \i+1\relax}
      \edef\ll{\number\numexpr \i\relax}
      \edef\l{\number\numexpr \i-1\relax}
      \coordinate
        (b\lll u)
        at ($ (b\ll u)!-1!(b\l u) $)
      ;
    \fi
  }}
  \coordinate (hs1) at ($ (h)+(0,5pt) $);
  \coordinate (hs2) at ($ (hs1)+(0,-12pt) $);

  \foreachbit \i/\bit in {0\orygbin}{
    \node[anchor=north]
      (p\i) at ($ (b\i u)+(hs2) $) {$\i$};
  }
  \node[anchor=east] (p) at (b.east |- p0) {pozycje:};

  \draw
    (b0u)
    -- (b\strlen{\orygbin}u)
    -- ($ (b\strlen{\orygbin}u)+(h) $)
    -- ($ (b0u)+(h) $)
    -- cycle
  ;
  \foreachbit \i/\bit in {0\orygbin}{
    \draw
      (b\i u) + (hs1)
      -- ++(hs2)
    ;
  }

  \draw[
    decorate,
    decoration={brace,mirror,amplitude=8pt},
  ]
    ($ (b\strlen{\orygbin}u)+(-3pt,5pt) $)
    -- ($ (b0u)+(3pt,5pt) $)
    node[midway,above=8pt] (n) {$\strlen{\orygbin}$}
  ;
  \node[anchor=east]
    (nd)
    at (b.east |- n)
    {długość kodu:}
  ;
\end{tikzpicture}

  \caption{
    Zobrazowanie oznaczenia pozycji występujących
    w~kodzie genetycznym przy krzyżowaniu,
    dla przykładowego $5$-cio znakowego
    binarnego genotypu: $01011$.
    Pozycję przed kodem binarnym oznacza $0$,
    każda następna wartość $i$ wskazuje pozycję
    bezpośrednio za $i$-tym elementem.
  }
  \label{fig:bin_gen}
\end{figure}

\subsection{Krzyżowanie -- część sekwencyjna}
\label{subsec:crossover_seq}
W~przypadku części sekwencyjnej zostanie
użyte krzyżowanie jednopunktowe,
gdzie losowo zostanie wybrany punkt podziału
i~części kodu genetycznego zostaną wymienione
między dwoma rodzicami tworząc parę dzieci.
Podgenotyp, posiadając stałą długość,
umożliwia prosty podział w~wybranym punkcie
i~wymianę fragmentów pomiędzy rodzicami.
Zachowuje to poprawność struktury chromosomu
oraz pozwala na efektywne łączenie cech
odziedziczonych po obu podosobnikach
przy niewielkiej złożoności obliczeniowej.

Przyjmując, że:\\
$n$ to długość genotypu podosobnika (liczba bitów),\\
$R\left(p, k\right)$ jest funkcją losową
  o~rozkładzie jednostajnym dyskretnym
  na zbiorze $\mathbb{Z} \cap \left[p, k\right]$,\\
$r_1\left[i\right]$ i~$r_2\left[i\right]$
  są rodzicami 1 i~2 dla
  $i \in \mathbb{N} \cap \left[0, n-1\right]$,\\
a~$d_1\left[i\right]$ i~$d_2\left[i\right]$
  są dziećmi 1 i~2 dla
  $i \in \mathbb{N} \cap \left[0, n-1\right]$,\\
to krzyżowanie jednopunktowe można
matematycznie opisać w~następujący sposób:
\vspace{-\parskip}
\begin{fleqn}
  \begin{gather}
    \label{math:seq_cp_rand}
    p_r=R\left(0, n\right)\\
    \label{math:seq_child1}
    d_1\left[i\right]=\begin{cases}
      r_1[i]
        ,& i < p_r\\
      r_2[i]
        ,& i \geqslant p_r\\
    \end{cases}\\
    d_2\left[i\right]=\begin{cases}
      r_2[i]
        ,& i < p_r\\
      r_1[i]
        ,& i \geqslant p_r\\
    \end{cases}
    \label{math:seq_child2}
  \end{gather}
\end{fleqn}

\Cref{math:seq_cp_rand} pokazuje losowanie
punktu krzyżowania z~dyskretnego zbioru
liczb naturalnych niewiększych niż $n$.
Oznacza to, że może być ono w~każdym punkcie podgenotypu,
nawet przed lub za (wylosowane odpowiednio $0$ lub $n$).
Jest ona pierwszą podoperacją
i~wartością zwrotną jest $p_r$.
\Cref{math:seq_child1,math:seq_child2}
pokazują działanie drugiej podoperacji,
która mając punkt krzyżowania tworzy dwóch potomków.

\subsection{Krzyżowanie -- część listowa}
\label{subsec:crossover_list}
W~przypadku części listowej zostanie
użyte krzyżowanie przeznaczone
do zmiennej długości genotypu,
inspirowane metodą przedstawioną
w~\citebook[cite={rys. 1}]{goldberg1989messy},
w~której przedstawione jest krzyżowanie dwupunktowe.
Natomiast jest on inny od opisanego
w~wspomnianym dokumencie algorytmu,
ponieważ o~ile pierwszy punkt podziału
jest wybierany z~pełnego zakresu
(biorąc pod uwagę wszystkie bity listy),
to punkt podziału drugiej listy jest wybierany
na podstawie wylosowanego elementu listy,
by zachować spójność długości list oraz
by zapobiec występowaniu ``półelementów''.
Dodatkowo algorytm pilnuje by lista
nie miała mniej elementów, niż minimalna liczba,
a~regulowanie do maksymalnej wartości jest
wykonywane przez przycinanie listy po krzyżowaniu.

Oznacza to, że pierwsze losowanie,
poza decydowaniem konkretnego przecięcia,
decyduje o~tym czy i~w~którym punkcie
nastąpi podział elementu w~drugiej liście.

Przyjmując, iż:\\
$e_s$ to wielkość pojedynczego elementu,\\
$l_{\min}$ i~$l_{\max}$ są minimalną
  i~maksymalną liczbą  elementów,\\
$R\left(p, k\right)$ jest funkcją losową
  o~rozkładzie jednostajnym dyskretnym na zbiorze
  $\mathbb{Z} \cap \left[p, k\right]$,\\
$n_{r_1}$ i~$n_{r_2}$ są liczbami
  elementów dla rodziców 1 i~2,\\
$r_1\left[i_{r_1}\right]$ i~$r_2\left[i_{r_2}\right]$ to
  rodzice 1 i~2 odpowiednio dla $
    i_{r_1} \in \mathbb{N}
    \cap \left[0, e_s n_{r_1}-1\right]
  $ i~$
    i_{r_2} \in \mathbb{N}
    \cap \left[0, e_s n_{r_2}-1\right]
  $,\\
$n_{d_1}$ i~$n_{d_2}$ są liczbami
  elementów dla dzieci 1 i~2,\\
$d_1\left[i_{d_1}\right]$ i~$d_2\left[i_{d_2}\right]$ to
  dzieci 1 i~2 odpowiednio dla $
    i_{d_1} \in \mathbb{N}
    \cap \left[0, e_s n_{d_1}-1\right]
  $ i~$
    i_{d_2} \in \mathbb{N}
    \cap \left[0, e_s n_{d_2}-1\right]
  $,\\
to krzyżowanie dwupunktowe będzie matematycznie
opisane jak w~poniższych podrozdziałach.
Pierwsza podoperacja zwraca typ
\texttt{tuple{\textless}int, int{\textgreater}}.
Opis jej został jednakże rozbity na
\cref{subsubsec:list_cp1,subsubsec:list_cp2},
gdzie każda z~nich zwraca po jednej
z~wartości \texttt{int}.
Druga podoperacja jest opisana
w~\cref{subsubsec:list_d1_d2}.

\subsubsection{Punkt przecięcia pierwszego rodzica}
\label{subsubsec:list_cp1}
\vspace{-\parskip}
\begin{fleqn}
  \begin{gather}
    \label{math:list_cp1_rand}
    p_{r_1}=R\left(0, e_s n_{r_1}\right)\\
    \label{math:list_b}
    l_b=p_{r_1} \mod e_s\\
    \label{math:list_calc1}
    \begin{gathered}
      l_{p_b}=\begin{cases}
        1 ,& l_b>0\\
        0 ,& l_b=0
      \end{cases}\\
      l_{p_{d_1}}=\left\lfloor
        \frac{p_{r_1}}{e_s}
      \right\rfloor + l_{p_b}\\
      l_{k_{d_2}}=n_{r_1}-\left\lfloor
        \frac{p_{r_1}}{e_s}
      \right\rfloor
    \end{gathered}
  \end{gather}
\end{fleqn}
\Cref{math:list_cp1_rand} pokazuje sposób
wylosowania punkt krzyżowania na pierwszym rodzicu.
Daje to możliwość obliczenia przesunięcia cięcia
na drugim rodzicu w~\cref{math:list_b}.
Następnie w~\cref{math:list_calc1} wykonywane
są obliczenia czy będzie przesunięcie ($l_{p_b}$)
oraz jaka jest gwarantowana minimalna długość pierwszego
potomka ($l_{p_{d_1}}$) i~drugiego ($l_{k_{d_2}}$)
na podstawie pierwszego wylosowanego punktu przecięcia.
Wartością zwrotną jest $p_{r_1}$, czyli punkt
krzyżowania na pierwszym z~rodziców.

\subsubsection{Punkt przecięcia drugiego rodzica}
\label{subsubsec:list_cp2}
\vspace{-\parskip}
\begin{fleqn}
  \begin{gather}
    \label{math:list_cp2_p}
    p_{r_{2_p}}=\begin{cases}
      0
        ,& l_{k_{d_2}} \geqslant l_{\min}\\
      l_{\min}-l_{k_{d_2}}
        ,& l_{k_{d_2}}<l_{\min}\\
    \end{cases}\\
    \label{math:list_cp2_k}
    p_{r_{2_k}}=\begin{cases}
      n_{r_2}-l_{p_b}
        ,& l_{p_{d_1}} \geqslant l_{\min}\\
      n_{r_2}-\left(l_{\min}-l_{p_{d_1}}+l_{p_b}\right)
        ,& l_{p_{d_1}}<l_{\min}\\
    \end{cases}\\
    \label{math:list_cp2_rand}
    p_{r_2}=e_s R\left(p_{r_{2_p}}, p_{r_{2_k}}\right)
      +l_b
  \end{gather}
\end{fleqn}
\Cref{math:list_cp2_p,math:list_cp2_k} obliczają punkt
początkowy losowania i~końcowy, zapobiega to przypadkowi,
iż któregoś z~potomków kod genetyczny
będzie za krótki względem wymagań.
Samo losowanie punktu krzyżowania na drugim rodzicu
jest wykonywane poprzez \cref{math:list_cp2_rand}.
Wynik tego równania jest wartością zwrotną.

\subsubsection{Powstanie kodu genetycznego potomków}
\label{subsubsec:list_d1_d2}
\vspace{-\parskip}
\begin{fleqn}
  \begin{gather}
    \label{math:list_delta}
    {\Delta p_r}=p_{r_1}-p_{r_2}\\
    \label{math:list_nd}
    \begin{gathered}
      n_{d_1}=\begin{cases}
        \frac{p_{r_1}
          + \left(e_s n_{p_2}-p_{r_2}\right)}{e_s}
          =n_{p_2}
            + \frac{\left(p_{r_1}-p_{r_2}\right)}{e_s}
          =n_{p_2}+\frac{{\Delta p_r}}{e_s}
            ,& n_{p_2}+\frac{{\Delta p_r}}{e_s}
              \leqslant l_{\max}\\
        l_{\max}
          ,& n_{p_2}+\frac{{\Delta p_r}}{e_s}>l_{\max}\\
      \end{cases}\\
      n_{d_2}=\begin{cases}
        \frac{p_{r_2}
          + \left(e_s n_{p_1}-p_{r_1}\right)}{e_s}
          =n_{p_1}
            - \frac{\left(p_{r_1}-p_{r_2}\right)}{e_s}
          =n_{p_1}-\frac{{\Delta p_r}}{e_s}
            ,& n_{p_1}-\frac{{\Delta p_r}}{e_s}
              \leqslant l_{\max}\\
        l_{\max}
          ,& n_{p_1}-\frac{{\Delta p_r}}{e_s}>l_{\max}\\
      \end{cases}
    \end{gathered}\\
    \begin{gathered}
      d_1\left[i_{d_1}\right]=\begin{cases}
        p_1[i_{d_1}]
          ,& i_{d_1} < p_{r_1}\\
        p_2[i_{d_1}-p_{r_1}+p_{r_2}]
          =p_2[i_{d_1}-{\Delta p_r}]
          ,& i_{d_1} \geqslant p_{r_1}\\
      \end{cases}\\
      d_2\left[i_{d_2}\right]=\begin{cases}
        p_2[i_{d_2}]
          ,& i_{d_2} < p_{r_2}\\
        p_1[i_{d_2}-p_{r_2}+p_{r_1}]
          =p_1[i_{d_2}+{\Delta p_r}]
          ,& i_{d_2} \geqslant p_{r_2}\\
      \end{cases}
    \end{gathered}
    \label{math:list_d}
  \end{gather}
\end{fleqn}
Powyżej widać wyliczenie pomocniczej
wartości przez \cref{math:list_delta}.
Jest ona używana przy obu
następnych zestawach obliczeń,
czyli obliczania ilości elementów potomków
w~\cref{math:list_nd} oraz bezpośrednio
ich postaci w~\cref{math:list_d}.
Dzielenie widoczne w~\cref{math:list_nd}
zawsze daje liczby całkowite,
co można udowodnić po przekształceniu
\cref{math:list_b} i~wprowadzeniu do dzielenia,
co jest pokazane poniżej:
\begin{equation}
  \displaystyle
  \label{math:list_div_prove}
  \begin{gathered}
    \displaystyle
    l_b=p_{r_1} \mod e_s
      \implies \exists {
        k: k \in \mathbb{Z}
        \land p_{r_1}=e_s k + l_b
      }\\
    \Downarrow\\
    \begin{aligned}
      \frac{{\Delta p_r}}{e_s}
        &=\frac{p_{r_1} - p_{r_2}}{e_s}
        =\frac{
          \left(e_s k + l_b\right)
          - \left(
            e_s R\left(p_{r_{2_p}}, p_{r_{2_k}}\right)
            + l_b
          \right)
        }{e_s}\\
        &=\frac{
          e_s\left(
            k - R\left(p_{r_{2_p}}, p_{r_{2_k}}\right)
          \right)
        }{e_s}
        =k - R\left(p_{r_{2_p}}, p_{r_{2_k}}\right)
    \end{aligned}
  \end{gathered}
\end{equation}
Na końcu przekształceń jest
$k - R\left(p_{r_{2_p}}, p_{r_{2_k}}\right)$,
oznacza to, że wynikiem dzielenia
jest zawsze liczba całkowita,
ponieważ $k \in \mathbb{Z}$
oraz wynik funkcji $R\left(p, k\right)$
jest również z~zakresu liczb całkowitych.

\subsection{Krzyżowanie -- osobnik}
\label{subsec:crossover_ind}
Krzyżowanie na poziomie osobnika polega
na delegowaniu zadań do podosobników.
Podczas pierwszej podoperacji,
polegającej na określeniu punktów krzyżowania,
osobnik odpytuje podosobnika sekwencyjnego
i~dwóch podosobników listowych łącząc je
za pomocą typu \texttt{tuple},
co zostało ukazane na \cref{fig:gen_parts}.
Na przytoczonym rysunku widać,
że typ punktu krzyżowania
dla podosobnika sekwencyjnego
to \texttt{int}, a~dla listowego
-- \texttt{tuple{\textless}int, int{\textgreater}}.
Po połączeniu wyników z~odpytań wychodzi
typ \texttt{tuple{\textless}%
  int,
  tuple{\textless}int, int{\textgreater},
  tuple{\textless}int, int{\textgreater}%
{\textgreater}}.

Część drugiej podoperacji przebiega podobnie,
klasa osobnika deleguje krzyżowanie do podosobników.
Następnie klasa łączy skrzyżowane osobniki tworząc
nowe obiekty i~dba by fenotypy podosobników były zgodne,
co znaczy, że długość list powinna być równa
dla listy ilości neuronów oraz o~jeden większa
dla listy typów aktywacji w~stosunku do parametru
liczby warstw ukrytych z~podosobnika sekwencyjnego.
Jeżeli lista jest zbyt długa, to
nadmiarowe elementy są usuwane z~fenotypu.
W~przypadku gdy jest zbyt krótka są
używane wartości ziaren generatora
do wygenerowania dalszych wartości podgenotypu.
Przyjmując, że:\\
$R\left(z,n,p,k\right)$ jest funkcją
  pseudolosową wykorzystującą generator liczb
  pseudolosowych zainicjalizowany ziarnem $z$;
  zwraca $n$-elementowy ciąg liczb całkowitych
  należących do dyskretnego przedziału
  $\mathcal{D}=\mathbb{Z} \cap \left[p, k\right]$;
  każdy element jest losowany niezależnie
  z~zbioru $\mathcal{D}$;
  zastosowanie tego samego ziarna gwarantuje otrzymanie
  tego samego ciągu dla tych samych parametrów;
  jeżeli jedynym różnym parametrem będzie $n$
  przykładowo $n_1$ i~$n_2$, gdzie $n_1>n_2$,
  to pierwsze $n_2$ elementów jest identyczne,
  czyli ciąg zwrócony dla $n_2$ stanowi prefiks
  ciągu zwróconego dla $n_1$,\\
$z$ to parametr ziarna generatora odpowiadający
  wydłużanemu podosobnikowi listowemu,\\
$n$ i~$n_o$ są liczbami elementów fenotypów
  odpowiednio aktualnego i~oczekiwanego,\\
$l\left[i\right]$ jest ciągiem wartości
  fenotypu podosobnika listowego dla
  $i~\in \mathbb{Z} \cap \left[0, n-1\right]$,\\
$l_o\left[i_o\right]$ to ciąg wartości fenotypu
  podosobnika listowego o~oczekiwanej
  liczbie elementów ($n_o$) dla
  $i_o \in \mathbb{Z} \cap \left[0, n_o-1\right]$,\\
można opisać całość matematycznie w~następujący sposób:
\begin{equation}
  \label{eq:too_short_long}
  \begin{gathered}
    n_d=\min\left\{n, n_o\right\}\\
    l_o\left[i_o\right]=\begin{cases}
      l\left[i_o\right]
        ,& i_o < n_d\\
      {
        R\left(z, n_o-n_d, p, k\right)
      }{\left[i_o-n_d\right]}
        ,& i_o \geqslant n_d\\
    \end{cases}
  \end{gathered}
\end{equation}

\section{Mutacja}
\label{sec:mutate}
Mutacja będzie odbywała się
przez odwracanie losowych bitów.
Decyzja o~mutacji będzie
podejmowana względem osobnika,
jeżeli będzie poddawany mutacji,
to następnie wystąpi losowanie
w~których częściach genotypu wystąpi.
Sama operacja polega na wylosowaniu miejsca mutacji,
a~następnie zamienieniu bitu na przeciwny.
Losowanie powinno się odbyć poprzez funkcją losową
o~rozkładzie jednostajnym dyskretnym na zbiorze
$\mathbb{Z} \cap \left[0, n-1\right]$,
gdzie $n$ to długość podgenotypu.

\section{Ocena przystosowania}
\label{sec:fittness}
Ocena przystosowania osobnika będzie odbywała się poprzez
nauczenie określonej liczby razy sieci o~parametrach
i~z~parametrami wynikającymi z~fenotypu osobnika.
Liczba ta będzie określona przez liczbę zrównolegleń.
Obliczony błąd dokładności walidacyjnej zostaje
przekształcony na przystosowanie za pomocą funkcji
korygującej $f$ i~uśrednione.
Dla celów statystycznych przekształcane
jest również dopasowanie testowe.

Jeżeli:\\
$n$ jest liczbą warstw ukrytych,\\
$n_{\max}$ to maksymalna liczba warstw ukrytych\\
oraz $e$ jest błądem dopasowania
  z~przedziału $\mathbb{R}_{+} \cup \left\{0\right\}$,
  gdzie $0$ oznacza brak błędu,
  czyli idealne dopasowanie,\\
to domyślna postać funkcji $f$ wygląda następująco:
\begin{equation}
  \label{eq:correction_f}
  f\left(e, n\right)
    =\max{\left\{\frac{1}{
      1 + e \left(
        1
        + 0.2 \frac{n-1}{n_{\max}-1}
      \right)
    }, 0\right\}}
\end{equation}
Współczynnik $0.2$ został dobrany poprzez eksperymenty.
Jest on kompromisem między jakością
dopasowania, a~złożonością sieci.
Dzięki takiej wartości liczba warstw
wpływa na ocenę osobnika,
jednakże dużo ważniejsze pozostaje jego dopasowanie.
W~ten sposób algorytm preferuje prostsze
architektury jedynie w~sytuacjach
przy porównywalnej jakości działania,
a~nie niwelują bardziej złożonych modeli
jeśli zapewniają istotnie lepsze wyniki.

Aby zapewnić wysoki poziom możliwości
dopasowania działania biblioteki
do indywidualnych potrzeb,
podawany jest do funkcji korygującej cały osobnik.
Więc w~miejscu $n$ w~\cref{eq:correction_f}
jest w~bibliotece osobnik z~którego można
uzyskać wszystkie parametry sieci,
jeżeli ktoś chce również, na przykład,
stosować kary dla liczby neuronów
w~poszczególnych warstwach.

Jeżeli użytkownik poda tylko
dane treningowe i~testowe, to:\\
Uczenie będzie odbywało się tylko na danych treningowych.
Po zakończeniu uczenia zostanie obliczony błąd
dopasowania na danych treningowych i~testowych,
z~czego to drugie zostanie zwrócone,
natomiast nie będzie używane w~procesie ewolucji,
a~jest w~celach statystycznych.
Zwrócona sieć będzie tą z~ostatniej epoki uczenia.

Jeżeli użytkownik poda dodatkowo dane walidacyjne, to:\\
Uczenie będzie odbywało się na danych treningowych,
natomiast po każdej epoce obliczana będzie wartość
funkcji straty na zbiorze walidacyjnym.
Jeżeli przez ustaloną liczbę kolejnych
epok nie nastąpi poprawa tej wartości,
proces uczenia zostanie zakończony
przed osiągnięciem maksymalnej liczby epok.
Zwrócona sieć będzie tą z~epoki,
która dała najmniejszą stratę na danych walidacyjnych.
Jeżeli najmniejsza wartość funkcji straty
będzie reprezentowana przez kilka epok,
to zostaną wybrane wagi tej z~najwcześniejszej epoki.

\section{Fenotyp}
\label{sec:fenotype}
Jak przedstawiono na \cref{fig:gen_parts}
genotyp osobnika podzielony jest na trzy części.
Poniżej zostanie przedstawiony przykład
przy którym postać genotypu podosobnika
\textbf{\subref{subfig:gen_parts_SI}}
ma następującą postać binarną: \texttt{%
  0001100010010110001010100000100110010101100101101001011%
}. A~listy liczb neuronów i~typów aktywacji,
które są zakodowane w~dwóch podosobnikach listowych
\textbf{\subref{subfig:gen_parts_LI}},
mają odpowiednio następujące genotypy:
\texttt{01000101101111110001}
oraz \texttt{00100100111001}.

Po przyjęciu że\\
$
  v\left(x\right)=\sum_{i=0}^{l-1}{x_i2^{l-i-1}}
$ jest funkcją zamieniającą kod binarny $
  x=\left(x_0,x_1,\dots,x_{l-1}\right)
  \land x_i \in \left\{0,1\right\}
$ w~wartość liczbową,
  gdzie $x_0$ oznacza najbardziej znaczący bit,\\
$v_{\max}$ i~$v_{\min}$ są odpowiednio
  maksymalną i~minimalną wartością
  poszukiwanej liczby należącej do zbioru
  liczb całkowitych ($\mathbb{Z}$),\\
$l$ to liczba naturalna ($\mathbb{N}$)
  oznaczająca długość kodu binarnego
  zamienianego na wartość liczbową,\\
$b=\left(b_0, b_1, \dots, b_{l-1}\right)$ jest
  kodem binarnym zamienianym na wartość
  liczbową gdzie $b_i \in \left\{0, 1\right\}$,\\
$m=\left(m_0, m_1, \dots, m_{l-1}\right)$ jest
  kodem binarnym wartości $v_{\max}-v_{\min}$,
  czyli największej wartości możliwej do uzyskania
  przed dodaniem przesunięcia wynikającego
  z~minimalnej wartości zakresu;
  prawdziwe jest więc dla niego równanie
  $v\left(m\right)=v_{\max}-v_{\min}$,\\
a~$x$ to wartość poszukiwanej wartości
  liczbowej kodu binarnego $b$,\\
każda wartość liczby z~sekwencji lub listy
obliczana jest w~następujący sposób:
% \vspace{-\parskip}
\begin{fleqn}
  \begin{gather}
    \label{math:first_greater_bin}
    m_b=\min{\left\{
      \left\{i| m_i > b_i\right\}
      \cup \left\{l\right\}
    \right\}}\\
    \label{math:indexes_of_greater_bins}
    b_m=\left\{i|b_i>m_i \lor i<m_b\right\}\\
    \label{math:corrected_bin}
    b'=\left(b'_0,b'_1,\dots,b'_{l-1}\right):
      b'_i=\begin{cases}
        0 ,& i \in b_m \\
        b_i ,& i \notin b_m \\
      \end{cases}\\
    \label{math:bin_value}
    x=v\left(b'\right)+v_{\min}
  \end{gather}
\end{fleqn}
\Cref{math:first_greater_bin} wskazuje pierwszy indeks
na którym bit kodu $m$ jest większy od $b$.
Jeżeli nie ma takiego bitu przyjmowane zostaje $l$.
Następnie w~\cref{math:indexes_of_greater_bins} do
$b_m$ wpisane są wszystkie indeksy mniejsze niż $m_b$
na których bity $b$ są większe od $m$.
Dzięki temu wartość poprawionego kodu $b'$
-- czyli $v\left(b'\right)$ -- z~\cref{math:corrected_bin}
jest na pewno mniejsza lub równa $v\left(m\right)$,
dodatkowo jeżeli $
  v\left(m\right) \geqslant v\left(b\right)
$ to $v\left(b\right)=v\left(b'\right)$.
Końcowa wartość $x$ jest obliczona
jak widać na \cref{math:bin_value}.
Liczba binarna zamieniana jest na dziesiętną
i~dodawana jest do niej minimalna
wartość danego parametru.
W~poniższych podrozdziałach będzie pokazana
zmiana bitów tylko po uprzednim wykazaniu,
iż wartość wynikowa jest większa niż maksymalna,
bez korekcji zapewnionej przez procedurę opisaną powyżej.

\subsection{Fenotyp -- część sekwencyjna}
\label{subsec:fenotype_seq}
Część sekwencyjna osobnika jest zależna
od \textbf{\subref{subfig:gen_parts_MNSI}},
a~następnie od \textbf{\subref{subfig:gen_parts_SI}}.
Ten drugi z~nich nie tworzy podfenotypu,
w~przeciwieństwie do pierwszego.
Ten po odczytaniu przechowywanego kodu binarnego
interpretuje go dzieląc na długości
odpowiadające optymalizowanym liczbom,
które są opisane w~\cref{sec:individual_genetic_code}.
Po rozdzieleniu podgenotypu kody są
przekształcane na wartości.
Poniżej znajduje się rozpisany przykładowy podgenotyp
\texttt{%
  0001100010010110001010100000100110010101100101101001011%
}.
\begin{nospitemize}
  \item liczba warstw ukrytych
    -- 5 bitów $\left[1, 32\right]$
    -- \texttt{00011}$\displaystyle
      {}^{\left(b\right)}
        =3^{\left(d\right)}+1
        =4
    $,
  \item ziarno generatora liczby neuronów
    -- 10 bitów $\left[0, 1023\right]$
    -- \texttt{0001001011}$\displaystyle
      {}^{\left(b\right)}
        =75^{\left(d\right)}+0
        =75
    $,
  \item ziarno generatora typów aktywacji
    -- 10 bitów $\left[0, 1023\right]$
    -- \texttt{0001010100}$\displaystyle
      {}^{\left(b\right)}
        =84^{\left(d\right)}+0
        =84
    $,
  \item optymalizator
    -- 1 bit $\left[0, 1\right]$
    (%
      0 {\textrightarrow} SGD,
      1 {\textrightarrow} Adam%
    ) -- \texttt{0}$\displaystyle
      {}^{\left(b\right)}
        =0^{\left(d\right)}+0
        =0
    \Rightarrow$ SGD,
  \item wskaźnik uczenia się (learning rate)
    -- 12 bitów $
      {2^{-14}} {\cdot} {\left[1, 2^{12}\right]}
    $
    -- \texttt{001001100101}$\displaystyle
      {}^{\left(b\right)}
        =613^{\left(d\right)}+1
    $\break$\displaystyle
        =614
      \Rightarrow 2^{-14} \cdot 614
        = 0.037475586
    $,
  \item liczba epok
    -- 7 bitów $\left[1, 128\right]$
    -- \texttt{0110010}$\displaystyle
      {}^{\left(b\right)}
        =50^{\left(d\right)}+1
        =51
    $,
  \item wielkość partii
    -- 10 bitów $\left[1, 1024\right]$
    -- \texttt{1101001011}$\displaystyle
      {}^{\left(b\right)}
        =843^{\left(d\right)}+1
        =844
    $.
\end{nospitemize}

\subsection{Fenotyp -- część listowa -- liczby neuronów}
\label{subsec:fenotype_list_nnum}
Części listowe osobnika zależą od
\textbf{\subref{subfig:gen_parts_MNLI}},
a~następnie od \textbf{\subref{subfig:gen_parts_LI}}.
Podobnie jak w~przypadku sekwencyjnej
części druga z~nich nie tworzy fenotypu.
Natomiast fenotyp części listowej jest
też modyfikowany bezpośrednio przez osobnika
w~celu dostosowania długości powstałych list
do liczby warstw ukrytych z~części sekwencyjnej.

\tocless\subsubsection{%
  Budowanie fenotypu
  (\getsubrefname{subfig:gen_parts_MNLI})%
}
Kod binarny zostaje tu podzielony
na części o~równych długościach.
Następnie każdy wydzielony kod przechodzi
przez proces przekształcania.
Poniżej został wydzielony na elementy
przykładowy podgenotyp \texttt{01000101101111110001}:
\begin{nospitemize}
  \item pierwszy element
    -- \texttt{0100010110}$\displaystyle
      {}^{\left(b\right)}
        =278^{\left(d\right)}+2
        =280
    $
  \item drugi element
    -- \texttt{1111110001}$\displaystyle
      {}^{\left(b\right)}
        =1009^{\left(d\right)}+2
        =1011 > 724
    $
\end{nospitemize}

\begin{figure}
  \centering
  % bin: 1111110001
% max: 1011010010
\begin{tikzpicture}
  \def\orygbin{1111110001}
  \def\maxbin{1011010010}
  \def\corrbin{1011010001}
  \ExplSyntaxOn
    \def\foreachbit#1/#2 in #3#4{%
      \int_zero:N \l_tmpa_int
      \exp_args:Ne \tl_map_inline:nn {#3}
        {
          \edef#1{\int_use:N \l_tmpa_int}
          \edef#2{##1}
          #4
          \int_incr:N \l_tmpa_int
        }
    }
    \NewDocumentCommand{\strlen}{m}
    {
      \exp_args:Ne \tl_count:n {#1}
    }
  \ExplSyntaxOff

  \foreachbit \i/\bit in \orygbin{
    \node (b\i) at (\i em,0) {$\bit$};
  }
  \node[left=.3cm of b0] (b) {kod binarny do poprawienia:};
  \coordinate
    (h)
    at ($ (b0.south west)-(b0.north west) $)
  ;
  \coordinate
    (b1u)
    at ($ (b0.north east)!.5!(b1.north west) $)
  ;
  \coordinate
    (b2u)
    at ($ (b1.north east)!.5!(b2.north west) $)
  ;
  \coordinate
    (b0u)
    at ($ (b1u)!-1!(b2u) $)
  ;
  \foreachbit \i/\bit in {\orygbin}{{
    \ifnum\i<2
    \else
      \edef\lll{\number\numexpr \i+1\relax}
      \edef\ll{\number\numexpr \i\relax}
      \edef\l{\number\numexpr \i-1\relax}
      \coordinate
        (b\lll u)
        at ($ (b\ll u)!-1!(b\l u) $)
      ;
    \fi
  }}
  \coordinate (hp) at ($ (h)+(0,-7pt) $);

  \foreachbit \i/\bit in {\orygbin}{
    \edef\ii{\number\numexpr \i+1\relax}
    \node[anchor=north]
      (p\i)
      at ($ (b\i u)!.5!(b\ii u)+(hp) $)
      {$\i$}
    ;
    \draw
      ($ (b\i u)!.5!(b\ii u) $) + (h)
      -- ++(hp)
    ;
  }
  \node[anchor=east] (p) at (b.east |- p0) {pozycje:};

  \draw
    (b0u)
    -- (b\strlen{\orygbin}u)
    -- ($ (b\strlen{\orygbin}u)+(h) $)
    -- ($ (b0u)+(h) $)
    -- cycle
  ;
  \foreachbit \i/\bit in {\orygbin}{
    \ifnum\i<1
    \else
      \draw
        (b\i u)
        -- ++(h)
      ;
    \fi
  }

  \foreachbit \i/\bit in \maxbin{
    \node (m\i) at (\i em,2em) {$\bit$};
  }
  \node[left=.3cm of m0]
    (m)
    {maksymalna postać kodu binarnego:}
  ;

  \coordinate
    (dbm)
    at ($ (m0.north west)-(b0.north west) $)
  ;
  \foreachbit \i/\bit in {0\orygbin}{
    \coordinate
      (m\i u)
      at ($ (b\i u)+(dbm) $)
    ;
  }
  \draw
    (m0u)
    -- (m\strlen{\orygbin}u)
    -- ($ (m\strlen{\orygbin}u)+(h) $)
    -- ($ (m0u)+(h) $)
    -- cycle
  ;
  \foreachbit \i/\bit in {\orygbin}{
    \ifnum\i<1
    \else
      \draw
        (m\i u)
        -- ++(h)
      ;
    \fi
  }
  \node[anchor=north east]
    (m_b_node)
    at ($ (p.south east)+(0,0)!.4!(h) $)
    {$m_b$:}
  ;
  \node[anchor=north east]
    (b_m_node)
    at ($ (m_b_node.south east)+(0,0)!.4!(h) $)
    {$b_m$:}
  ;
  \node[right=0 of m_b_node,yshift=1pt]
    (m_b)
    {$8$}
  ;
  \node[right=0 of b_m_node]
    (b_m)
    {$\left\{1, 4\right\}$}
  ;












  \foreachbit \i/\bit in \corrbin{
    \node (corr\i) at (\i em,-8em) {$\bit$};
  }
  \node[left=.3cm of corr0] (corr) {poprawiony kod binarny:};

  \coordinate
    (dbc)
    at ($ (corr0.north west)-(b0.north west) $)
  ;
  \foreachbit \i/\bit in {0\orygbin}{
    \coordinate
      (corr\i u)
      at ($ (b\i u)+(dbc) $)
    ;
  }
  \draw
    (corr0u)
    -- (corr\strlen{\orygbin}u)
    -- ($ (corr\strlen{\orygbin}u)+(h) $)
    -- ($ (corr0u)+(h) $)
    -- cycle
  ;
  \foreachbit \i/\bit in {\orygbin}{
    \ifnum\i<1
    \else
      \draw
        (corr\i u)
        -- ++(h)
      ;
    \fi
  }

  \fill[blue!75!yellow, opacity=.3]
    (m8u)
    -- (m9u)
    -- ($ (m9u)+(h) $)
    -- ($ (m9u)!0.23!(m8u)+(h) $)
    -- ($ (b9u)!0.23!(b8u) $)
    -- (b9u)
    -- ($ (b9u)+(h) $)
    -- ($ (b8u)+(h) $)
    -- (b8u)
    -- ($ (b8u)!0.23!(b9u) $)
    -- ($ (m8u)!0.23!(m9u)+(h) $)
    -- ($ (m8u)+(h) $)
    -- cycle
  ;

  \fill[red, opacity=.3]
    (m1u)
    -- (m2u)
    -- ($ (m2u)+(h) $)
    -- ($ (m2u)!0.23!(m1u)+(h) $)
    -- ($ (b2u)!0.23!(b1u) $)
    -- (b2u)
    -- ($ (b2u)+(h) $)
    -- ($ (b1u)+(h) $)
    -- (b1u)
    -- ($ (b1u)!0.23!(b2u) $)
    -- ($ (m1u)!0.23!(m2u)+(h) $)
    -- ($ (m1u)+(h) $)
    -- cycle
  ;

  \fill[red, opacity=.3]
    (m4u)
    -- (m5u)
    -- ($ (m5u)+(h) $)
    -- ($ (m5u)!0.23!(m4u)+(h) $)
    -- ($ (b5u)!0.23!(b4u) $)
    -- (b5u)
    -- ($ (b5u)+(h) $)
    -- ($ (b4u)+(h) $)
    -- (b4u)
    -- ($ (b4u)!0.23!(b5u) $)
    -- ($ (m4u)!0.23!(m5u)+(h) $)
    -- ($ (m4u)+(h) $)
    -- cycle
  ;

  \fill[green!30, opacity=.3]
    (corr0u)
    -- (corr1u)
    -- ($ (corr1u)+(h) $)
    -- ($ (corr0u)+(h) $)
    -- cycle
  ;
  \fill[green!30, opacity=.3]
    (corr2u)
    -- (corr4u)
    -- ($ (corr4u)+(h) $)
    -- ($ (corr2u)+(h) $)
    -- cycle
  ;
  \fill[green!30, opacity=.3]
    (corr5u)
    -- (corr10u)
    -- ($ (corr10u)+(h) $)
    -- ($ (corr5u)+(h) $)
    -- cycle
  ;
  \fill[green, opacity=.3]
    (corr1u)
    -- (corr2u)
    -- ($ (corr2u)+(h) $)
    -- ($ (corr1u)+(h) $)
    -- cycle
  ;
  \fill[green, opacity=.3]
    (corr4u)
    -- (corr5u)
    -- ($ (corr5u)+(h) $)
    -- ($ (corr4u)+(h) $)
    -- cycle
  ;
\end{tikzpicture}

  \caption{
    Korekcja kodu binarnego względem
    jego maksymalnej dopuszczalnej postacji.
    Wartość $m_b$ jest pierwszą pozycją,
    w~której maksymalna postać kodu binarnego
    ma bit o~większej wartości niż korygowany.
    Bit ten został pokazany kolorem
    fioletowym~(\showcolor{blue!75!yellow}{.3}).
    Następnie do $b_m$ są wpisane indeksy mniejsze
    niż $m_b$ na których to korygowany kod
    jest większy niż maksymalna postać.
    Przypadki te zostały oznaczone
    kolorem czerwonym~(\showcolor{red}{.3}).
    Skorygowany kod został oznaczony
    kolorem zielonym~(\showcolor{green!30}{.3})
    z~szczególnym uwzględnieniem zmienionych
    bitów ciemniejszym kolorem~(\showcolor{green}{.3})
  }
  \label{fig:bin_corr}
\end{figure}  

Drugi element jest większy niż dopuszczalna wartość,
więc \texttt{1111110001} zostaje
skorygowane względem \texttt{1011010010},
które jest binarną reprezentacją wartości
$m=v_{\max}-v_{\min}=724-2=722$.
Korekcja wystąpi w~sposób przedstawiony
na \cref{fig:bin_corr}.
Drugi element po korekcie będzie miał wartość:
\texttt{1011010001}$\displaystyle
  {}^{\left(b\right)}
    =721^{\left(d\right)}+2
    =723
$

\tocless\subsubsection{%
  Korekta fenotypu
  (\getsubrefname{subfig:gen_parts_NNI})%
}
Osobnik przetwarza fenotyp pod względem długości.
Jeżeli jest on za długi, to zostają
wycięte nadmiarowe elementy.
W~przypadku gdy jest za krótki to jest uzupełniany.
Proces ten opisany jest w~\cref{subsec:crossover_ind}.
Poniżej został przedstawiony przykład działania
\cref{eq:too_short_long} dla otrzymanego fenotypu
$\left(280, 722\right)$, ziarna generatora $75$
oraz oczekiwanej długości $4$.
$$
  \begin{aligned}
    n_o-n_d&=4-2=2\\
    l_o&=\left(
      280,
      722,
      {R\left(75, 2, 2, 724\right)}{\left[0\right]},
      {R\left(75, 2, 2, 724\right)}{\left[1\right]}
    \right)\\
      &=\left(280, 722, 463, 600\right)
  \end{aligned}
$$

\subsection{Fenotyp -- część listowa -- typy aktywacji}
\label{subsec:fenotype_list_tactiv}
Część listowa typu aktywacji podlega
tym samym zależnością co opisane
w~\cref{subsec:fenotype_list_nnum}.
Różnice występują tu w~parametrach
własnych tej części listowej,
czyli w~jej dopuszczalnych wartościach
minimalnych i~maksymalnych oraz długości listy,
która będzie zawsze korygowana do wartości
o~jeden większej niż liczba warstw ukrytych.

\tocless\subsubsection{%
  Budowanie fenotypu
  (\getsubrefname{subfig:gen_parts_MNLI})%
}
Kod binarny zostaje podzielony
na części o~równych długościach.
Następnie każdy wydzielony kod przechodzi
przez proces przekształcania.
Poniżej został wydzielony na elementy
przykładowy podgenotyp \texttt{00100100111001}:
\begin{nospitemize}
  \item pierwszy element
    -- \texttt{00}$\displaystyle
      {}^{\left(b\right)}
        =0^{\left(d\right)}+0
        =0
    \Rightarrow$ relu
  \item drugi element
    -- \texttt{10}$\displaystyle
      {}^{\left(b\right)}
        =2^{\left(d\right)}+0
        =2
    \Rightarrow$ sigmoid
  \item trzeci element
    -- \texttt{01}$\displaystyle
      {}^{\left(b\right)}
        =1^{\left(d\right)}+0
        =1
    \Rightarrow$ tanh
  \item czwarty element
    -- \texttt{00}$\displaystyle
      {}^{\left(b\right)}
        =0^{\left(d\right)}+0
        =0
    \Rightarrow$ relu
  \item piąty element
    -- \texttt{11}$\displaystyle
      {}^{\left(b\right)}
        =3^{\left(d\right)}+0
        =3
    \Rightarrow$ leaky relu
  \item szósty element
    -- \texttt{10}$\displaystyle
      {}^{\left(b\right)}
        =2^{\left(d\right)}+0
        =2
    \Rightarrow$ sigmoid
  \item siódmy element
    -- \texttt{01}$\displaystyle
      {}^{\left(b\right)}
        =1^{\left(d\right)}+0
        =1
    \Rightarrow$ tanh
\end{nospitemize}

\tocless\subsubsection{%
  Korekta fenotypu
  (\getsubrefname{subfig:gen_parts_NNI})%
}
Osobnik przetwarza fenotyp pod względem długości.
Jeżeli jest on za długi, to wycina nadmiarowe elementy.
W~przypadku gdy jest za krótki to
jest uzupełniany do zadanej długości.
Proces ten został opisany w~\cref{subsec:crossover_ind}.
Poniżej przedstawiony jest przykład działania
\cref{eq:too_short_long} dla otrzymanego fenotypu $\left(
  \text{relu},
  \text{sigmoid},
  \text{tanh},
  \text{relu},
  \text{leaky relu},
  \text{sigmoid},
  \text{tanh}
\right)$, ziarna generatora $84$
oraz oczekiwanej długości $4+1=5$.
$$\begin{aligned}
  n_o-n_d&=5-7=-2\\
  l&=\left(
    \text{relu},
    \text{sigmoid},
    \text{tanh},
    \text{relu},
    \text{leaky relu},
    \text{sigmoid},
    \text{tanh}
  \right)\\
  l_o&=\left(
    \text{relu},
    \text{sigmoid},
    \text{tanh},
    \text{relu},
    \text{leaky relu}
  \right)\\
\end{aligned}$$

\subsection{Fenotyp -- osobnik}
\label{subsec:fenotype_ind}
Osobnik po skonstruowaniu pełnego fenotypu
podosobników reprezentuje model nienauczonej sieci.
Dopiero podczas oceny przystosowania
model jest faktycznie konstruowany.
Jednakże jeszcze przed badaniem wiadomo jaka będzie
architektura dzięki jednoznaczności reprezentacji.
Poniżej została przedstawiona konstrukcja
dla przykładowych fenotypów podosobników
z~\cref{%
  subsec:fenotype_seq,%
  subsec:fenotype_list_nnum,%
  subsec:fenotype_list_tactiv%
}.

Fenotyp podosobnika widocznego na
\subrefparens{fig:gen_parts}{subfig:gen_parts_MNSI}:
\begin{nospitemize}
  \item liczba warstw ukrytych -- $4$,
  \item ziarno generatora liczby neuronów -- $75$,
  \item ziarno generatora typów aktywacji -- $84$,
  \item optymalizator -- SGD,
  \item wskaźnik uczenia się (learning rate)
    -- $0.074951171875$,
  \item liczba epok -- $51$,
  \item wielkość partii -- $844$.
\end{nospitemize}

Fenotypy podosoników typu pokazanego na
\subrefparens{fig:gen_parts}{subfig:gen_parts_MNLI}:
\begin{nospitemize}
  \item liczby neuronów
    -- $\left(280, 722, 463, 600\right)$,
  \item typów aktywacji -- $\left(
    \text{relu},
    \text{sigmoid},
    \text{tanh},
    \text{relu},
    \text{leaky relu}
  \right)$.
\end{nospitemize}

\begin{figure}
  \centering \begin{tikzpicture}[
    >=Latex,
    neuron/.style={
        circle,
        draw,
        thick,
        minimum size=7mm,
        fill=blue!10,
    },
    neuron_dots/.style={
        font=\LARGE\bfseries,
    },
    neuron_num/.style={
        font=\large,
    },
    every node/.style={
      font=\small,
    }
]
  \makeatletter
    \DeclareRobustCommand
      \myvdots{\vbox{\baselineskip4\p@ \lineskiplimit\z@%
        \hbox{.}\vspace{3pt}\hbox{.}\vspace{3pt}\hbox{.}}}
  \makeatother

  % Input
  \node[
    neuron,
    fill=green!10,
  ] (I-n) at (0,0) {};
  \node[left=2mm of I-n.west] {$x_{n-1}$};
  \node[
    neuron_dots,
    above=2mm of I-n,
  ] (I-d) {\myvdots};

  \foreach \jprev/\j/\txt in {
    d/3/$x_3$,
    3/2/$x_2$,
    2/1/$x_1$,
    1/0/$x_0$
  } {
    \node[
      neuron,
      fill=green!10,
      above=2mm of I-\jprev,
    ] (I-\j) {};
    \node[left=2mm of I-\j.west] {\txt};
  }

  \node[
    neuron_num,
    above=4mm of I-0
  ] (n) {$\boldsymbol{n}$};

  % Hidden layers
  \coordinate (Ic0-1) at ($ (I-0)!.5!(I-1) $);
  \coordinate[below=-2mm of I-n] (Ib-n);
  \foreach \i/\num/\act in {
    1/{280}/{ReLU},
    2/{722}/{Sigmoid},
    3/{463}/{Tanh},
    4/{600}/{ReLU}
  } {{
    \edef\x{\number\numexpr 2*\i\relax}
    \node[neuron] (H-\i-0) at ($ (Ic0-1)+(\x,0) $) {};
    \foreach \j in {1,2} {
      \edef\jprev{\number\numexpr \j-1\relax}
      \node[
        neuron,
        below=2mm of H-\i-\jprev,
      ] (H-\i-\j) {};
    }
    \node[
      neuron_dots,
      below=2mm of H-\i-2,
    ] (H-\i-d) {\myvdots};
    \node[
      neuron,
      below=2mm of H-\i-d,
    ] (H-\i-n) {};

    \node[
      neuron_num,
      anchor=south,
    ]
      at (n.south-|H-\i-0)
      {$\boldsymbol{\num}$}
    ;
    \node[anchor=north] at (Ib-n-|H-\i-n) {\act};
  }}

  % Output
  \coordinate (H-4c0-1) at ($ (H-4-0)!.5!(H-4-1) $);
  \node[
    neuron,
    fill=red!10,
  ] (O-0) at ($ (H-4c0-1)+(2,0) $) {};
  \node[
    neuron,
    fill=red!10,
    below=2mm of O-0,
  ] (O-1) {};
  \node[
    neuron_dots,
    below=2mm of O-1,
  ] (O-d) {\myvdots};
  \node[
    neuron,
    fill=red!10,
    below=2mm of O-d,
  ] (O-n) {};
  \node[right=4mm of O-0.east] {$y_0$};
  \node[right=4mm of O-1.east] {$y_1$};
  \node[right=4mm of O-n.east] {$y_{m-1}$};
  \node[anchor=north] at (Ib-n-|O-n) {Leaky ReLU};
  \node[
    neuron_num,
    anchor=south,
  ]
    at (n.south-|O-0)
    {$\boldsymbol{m}$}
  ;


  % Input -> H1
  \foreach \i in {0,1,2,3,n}
    \foreach \j in {0,1,2,n}
      \draw[gray!70] (I-\i.east)--(H-1-\j.west);

  % Hn-1 -> Hn
  \foreach \L/\R in {1,2,3} {
    \edef\R{\number\numexpr \L+1\relax}
    \foreach \i in {0,1,2,n}
      \foreach \j in {0,1,2,n}
        \draw[gray!70] (H-\L-\i.east)--(H-\R-\j.west);
  }

  % H4 -> Output
  \foreach \i in {0,1,2,n}
    \foreach \j in {0,1,n}
      \draw[gray!70] (H-4-\i.east)--(O-\j.west);

  % Output -> y
  \foreach \i in {0,1,n}
    \draw[gray!70,->] (O-\i.east) -- ++(10pt, 0);

  % ---
  % Braces
  % ---
  \draw[
    decorate,
    decoration={brace,mirror,amplitude=8pt},
  ]
    ($ (I-n)!-0.5!(I-n-|H-1-0)+(0.1,-.9) $)
    --($ (I-n)!0.5!(I-n-|H-1-0)+(-0.1,-.9) $)
    node[midway,yshift=-0.9cm]
      {\shortstack{Warstwa\\wejściowa}}
  ;

  \draw[
    decorate,
    decoration={brace,mirror,amplitude=8pt},
  ]
    ($ (I-n)!0.5!(I-n-|H-1-0)+(0.1,-.9) $)
    --($ (I-n-|H-4-0)!0.5!(I-n-|O-0)+(-0.1,-.9) $)
    node[midway,yshift=-0.9cm]
      {\shortstack{Warstwy\\ukryte}}
  ;

  \draw[
    decorate,
    decoration={brace,mirror,amplitude=8pt},
  ]
    ($ (I-n-|O-0)!0.5!(I-n-|H-4-0)+(0.1,-.9) $)
    --($ (I-n-|O-0)!-0.5!(I-n-|H-4-0)+(-0.1,-.9) $)
    node[midway,yshift=-0.9cm]
      {\shortstack{Warstwa\\wyjściowa}}
  ;
\end{tikzpicture}

  \caption{
    Architektura sieci neuronowej
    dla przykładowych danych
    z~$n$~atrybutami wejściowymi
    oraz $m$~wyjściowymi.
    Model jest w~pełni połączonym
    wielowarstwowym perceptronem.
    Ma 4 warstwy o~różnej
    liczbie neuronów w~warstwach,
    kolejno $280$, $722$, $463$ i~$600$.
    Mają one różne funkcje aktywacji,
    odpowiednio ReLU, Sigmoid, Tanh oraz ReLU.
    Funkcję aktywacji posiada również
    warstwa wyjściowa, tu Leaky ReLU.
  }
  \label{fig:ind_fenotype}
\end{figure}

Powyższy fenotyp daje architekturę
przedstawioną na \cref{fig:ind_fenotype}.
W~przykładzie nie została wyspecyfikowana
ilość neuronów wejściowych ani wyściowych,
ponieważ nie zostały przyjęte żadne przykładowe dane.

\begin{table}
  \centering
  \begin{tabular}{|c|c|c|c|c|c|c|c|c|c|}
  \hline
    $x_1$ & $x_2$ & $x_3$ & $x_4$ & $x_5$ & $x_6$
    & $y_1$ & $y_2$ & $y_3$ & $y_4$ \\
  \hline
    $0.767$ & $0.352$ & $0.663$ & $0.876$ & $0.612$
    & $0.094$ & $0.287$ & $0.218$ & $0.626$ & $0.369$ \\
    $0.670$ & $0.886$ & $0.535$ & $0.031$ & $0.073$
    & $0.239$ & $0.519$ & $0.723$ & $0.639$ & $0.717$ \\
    \vdots & \vdots & \vdots & \vdots &\vdots
    & \vdots & \vdots & \vdots & \vdots & \vdots\\
    $0.898$ & $0.884$ & $0.735$ & $0.082$ & $0.453$
    & $0.533$ & $0.066$ & $0.891$ & $0.066$ & $0.173$ \\
  \hline
\end{tabular}

  \caption{
    Przykładowe dane, gdzie atrybuty wejściowe
    i~wyjściowe znajdują się w~jednej tabeli.
    Pierwszych jest 6 i~mają następujące nazwy:
    $x_1$, $x_2$, $x_3$, $x_4$, $x_5$ oraz $x_6$.
    A~4 atrybuty wyjściowe to:
    $y_1$, $y_2$, $y_3$ oraz $y_4$.
  }
  \label{tab:example_data}
\end{table}

Jednakże oznacza to również, że po uruchomieniu programu,
wiadome jest jakie są wartości $n$~i~$m$,
ponieważ ilość atrybutów wraz z~samymi
danymi są podawane przez użytkownika,
by program mógł uczyć sieci
i~obliczać ich przystosowanie.
Jeżeli dane miałyby postać
pokazaną w~\cref{tab:example_data},
a~do funkcji zostałoby podany parametr
liczby atrybutów wejściowych równy
$6$ lub $-4$, to $n=6 \land m=4$.
W~takim przypadku liczby neuronów kolejno
we wszystkich warstwach wyglądałyby następująco:
$6$, $280$, $722$, $463$, $600$ i~$4$.

\subsubsection*{Wyjątek}
Powyżej został opisany przykład konstrukcji
dla danych o~problematyce regresyjnej.
W~przypadku problematyki klasyfikacyjnej
warstwa wyjściowa ma zawsze funkcję
aktywacji typu \texttt{softmax},
a~ostatni element podfenotypu listowego
typów aktywacji jest ignorowany.

\section{Implementacja}
\label{sec:implementation}
Implementacja zostanie pokazana odpowiednio
do poruszanych wcześniej podrozdziałów.

\subsection{%
  \refbykey{sec:proj_struct}{name}%
}
Biblioteka udostępnia dwie funkcje \texttt{cr\_network}
i~\texttt{cr\_net\_from\_ind} wraz z~stałymi
\texttt{const} oraz klasy \texttt{Generations}
i~\texttt{NetIndividual}.
\inputminted[firstline=4]{python}{\codedir/__init__.py}

Główna funkcja ma nazwę \texttt{cr\_network}.
Wpierw przetwarza czy są podane dane walidacyjne
oraz zapewnia, że interface tabeli będzie jednolity
z~typem \texttt{ndarray}.
\inputminted[
  firstline=91,
  lastline=96,
  % autogobble,
]{python}{\codedir/core.py}

Następnie wywoływana jest funkcja \texttt{get\_fit\_func}
przygotowująca dane i~ich sprawdzanie.
\inputminted[
  firstline=98,
  lastline=105,
  % autogobble,
]{python}{\codedir/core.py}

Po powrocie sterowania do głównej
funkcji tworzy ona obiekt pokoleń
i~oddelegowuje przeprowadznie ewolucji.
\inputminted[
  firstline=141,
  lastline=142,
  % autogobble,
]{python}{\codedir/core.py}

Funkcja \texttt{get\_fit\_func}
sprawdza czy tabele są dwuwymiarowe.
\inputminted[
  firstline=203,
  lastline=208,
  % autogobble,
]{python}{\codedir/utils.py}

Następnie czy mają równe ilości atrybutów:
\inputminted[
  firstline=209,
  lastline=212,
  % autogobble,
]{python}{\codedir/utils.py}

Ostatecznie czy dane są typu
deklarowanego przy uruchomieniu.
Mogą być regresyjne lub klasyfikacyjne.
\inputminted[
  firstline=261,
  lastline=270,
  % autogobble,
]{python}{\codedir/utils.py}

Zwracana funkcja oceniająca
dopasowania walidacyjne oraz testowe,
jest definiowana tak by przyjmować tylko osobnika
oraz ścieżkę folderu do zapisu danych z~uczenia.
\inputminted[
  firstline=286,
  lastline=289,
  % autogobble,
]{python}{\codedir/utils.py}

Funkcja ta tworzy model.
\inputminted[
  firstline=290,
  lastline=295,
  % autogobble,
]{python}{\codedir/utils.py}

Uczy stworzoną sieć.
\inputminted[
  firstline=304,
  lastline=320,
  % autogobble,
]{python}{\codedir/utils.py}

I~oblicza błąd dopasowania walidacyjnego.
\inputminted[
  firstline=322,
  lastline=329,
  % autogobble,
]{python}{\codedir/utils.py}

Oraz liczony jest błąd dopasowania do danych testowych.
\inputminted[
  firstline=330,
  lastline=334,
  % autogobble,
]{python}{\codedir/utils.py}

Ostatecznie przed zwróceniem danych zapisane
zostają wagi nauczonego modelu.
\inputminted[
  firstline=344,
  lastline=344,
  % autogobble,
]{python}{\codedir/utils.py}

I~metadane: liczba epok, wielkość wsadu, backend,
liczba warstw ukrytych, historia strat podczas kolejnych
epok oraz błędy dopasowań walidacyjnych i~testowych.
\inputminted[
  firstline=345,
  lastline=357,
  % autogobble,
]{python}{\codedir/utils.py}

Zwrócone zostają obie wartości błędów dopasowania.
\inputminted[
  firstline=358,
  lastline=358,
  % autogobble,
]{python}{\codedir/utils.py}

\subsection{%
  \refbykey{sec:generations_class}{name}%
}
Obiekt pokoleń tworzy obiekt populacji.
\inputminted[
  firstline=19,
  lastline=19,
  % autogobble,
]{python}{\codedir/classes/Generations.py}

Po dostaniu polecenia przeprowadzenia ewolucji sprawdza
ile jeszcze zostało ich do wykonania.
\inputminted[
  firstline=35,
  lastline=38,
  % autogobble,
]{python}{\codedir/classes/Generations.py}

Jeżeli nie została osiągnięta zadana ilość
\inputminted[
  firstline=41,
  lastline=41,
  % gobble=4,
]{python}{\codedir/classes/Generations.py}
to wywoływana jest następna generacja
i~pobierane są następnie wyniki.
\inputminted[
  firstline=45,
  lastline=47,
  % gobble=4,
]{python}{\codedir/classes/Generations.py}

\subsection{%
  \refbykey{sec:population_class}{name}%
}
Obiekt populacji z~którym występuje komunikacja powyżej od
razu podczas jego tworzenia generuje populację początkową.
\inputminted[
  firstline=42,
  lastline=42,
  % autogobble,
]{python}{\codedir/classes/Population.py}

Generowanie osobników do populacji
początkowej składa się z~dwóch etapów.
Pierwszym etapem jest wygenerowanie części sekwencyjnej.
\inputminted[
  firstline=54,
  lastline=54,
  % autogobble,
]{python}{\codedir/classes/NetIndividual.py}

Które następuje przez
wylosowanie zadanej liczby bitów.
\inputminted[
  firstline=15,
  lastline=15,
  % autogobble,
]{python}{\codedir/classes/GenIndividual.py}

Drugi etap generowania rozpoczyna się odczytaniem
liczby ukrytych warstw dla danego osobnika
z~wygenerowanej części sekwencyjnej.
\inputminted[
  firstline=55,
  lastline=55,
  % autogobble,
]{python}{\codedir/classes/NetIndividual.py}

A~następnie są generowane części listowe.
\inputminted[
  firstline=56,
  lastline=69,
  % autogobble,
]{python}{\codedir/classes/NetIndividual.py}

Które odbywa się przez
wylosowanie zadanej liczby bitów.
\inputminted[
  firstline=25,
  lastline=26,
  % autogobble,
]{python}{\codedir/classes/GenIndividual.py}

Następnie obliczane jest przystosowanie
osobników należących do populacji początkowej.
\inputminted[
  firstline=62,
  lastline=62,
  % autogobble,
]{python}{\codedir/classes/Population.py}

Podczas liczenia przystosowania taka sama architektura
sieci jest poddawana uczeniu zadaną ilość razy.
\inputminted[
  % firstline=73,
  firstline=75,
  lastline=79,
  % autogobble,
]{python}{\codedir/classes/Population.py}

Populacja jest w~stanie wygenerować następną generację
po wywołaniu metody \texttt{next\_generation}.
Metoda ta wykonuje jedno pełne przejście generacji,
czyli selekcję, krzyżowanie, mutację
i~obliczenie przystosowania osobników.
\inputminted[
  firstline=167,
  lastline=176,
  % autogobble,
]{python}{\codedir/classes/Population.py}

\subsection{%
  \refbykey{sec:individual_genetic_code}{name}%
}
Kod genetyczny osobnika składa się z~trzech podosobników:
\inputminted[
  firstline=70,
  lastline=70,
  % autogobble,
]{python}{\codedir/classes/NetIndividual.py}

Parametry części sekwencyjnej:
\inputminted[
  firstline=12,
  lastline=21,
]{python}{\codedir/consts.py}

Parametry części listowej:
\inputminted[
  firstline=23,
  lastline=25,
]{python}{\codedir/consts.py}

Uwzględnienie jednego elementu
więcej dla listy typów aktywacji:
\inputminted[
  firstline=18,
  lastline=18,
  % autogobble,
]{python}{\codedir/classes/NetIndividual.py}

\subsection{%
  \refbykey{sec:selection}{name}%
}
Użycie selekcji losowej w~przypadku gdy wszystkie
wartości przystosowania są równe zero:
\inputminted[
  firstline=129,
  lastline=130,
  % autogobble,
]{python}{\codedir/classes/Population.py}

Żadna wartość nie jest dodana do przystosowań,
jeżeli żadne nie jest równe 0:
\inputminted[
  firstline=132,
  lastline=136,
  % autogobble,
]{python}{\codedir/classes/Population.py}

Obliczenie wartości do dodania do przystosowań:
\inputminted[
  firstline=121,
  lastline=126,
  % autogobble,
]{python}{\codedir/classes/Population.py}

Obliczenie skumulowanego prawdopodobieństwa
do ruletki z~możliwymi wartościami
$\left[0, 1\right]$ wybrania osobników:
\inputminted[
  firstline=138,
  lastline=143,
  % autogobble,
]{python}{\codedir/classes/Population.py}

Wybranie osobników za pomocą ruletki:
\inputminted[
  firstline=146,
  lastline=148,
  % autogobble,
]{python}{\codedir/classes/Population.py}

\subsection{%
  \refbykey{sec:crossover}{name}%
}
Krzyżowanie jest inaczej realizowane
przez dwa typy podosobników.

\subsubsection{%
  \refbykey{subsec:crossover_seq}{name}%
}
Losowanie punktu krzyżowania:
\inputminted[
  firstline=36,
  lastline=38,
  % autogobble,
]{python}{\codedir/classes/GenIndividual.py}

Początek krzyżowania:
\inputminted[
  firstline=41,
  lastline=42,
  % autogobble,
]{python}{\codedir/classes/GenIndividual.py}

I~dwa razy wykonuje się poniższe
z~zamienionymi \texttt{a}~i~\texttt{b}:
\inputminted[
  firstline=22,
  lastline=22,
  % autogobble,
]{python}{\codedir/classes/GenIndividual.py}

\subsubsection{%
  \refbykey{subsec:crossover_list}{name}%
}
Losowanie punktów krzyżowania:
\inputminted[
  firstline=66,
  lastline=81,
  % autogobble,
]{python}{\codedir/classes/ListIndividual.py}

Początek krzyżowania:
\inputminted[
  firstline=84,
  lastline=85,
  % autogobble,
]{python}{\codedir/classes/ListIndividual.py}

I~dwa razy wykonuje się poniższe
z~zamienionymi \texttt{a}~i~\texttt{b}
oraz zamienioną kolejnością w~\texttt{cross\_point}:
\inputminted[
  firstline=44,
  lastline=44,
  % autogobble,
]{python}{\codedir/classes/ListIndividual.py}

\subsubsection{%
  \refbykey{subsec:crossover_ind}{name}%
}
Losowanie punktów krzyżowania:
\inputminted[
  firstline=125,
  lastline=132,
  % autogobble,
]{python}{\codedir/classes/NetIndividual.py}

Początek krzyżowania:
\inputminted[
  firstline=135,
  lastline=140,
  % autogobble,
]{python}{\codedir/classes/NetIndividual.py}

I~dwa razy wykonuje się poniższe
z~zamienionymi \texttt{a}~i~\texttt{b}
oraz zamienioną kolejnością w~dwóch
ostatnich częściach \texttt{cross\_point}:
\inputminted[
  firstline=77,
  lastline=85,
  % autogobble,
]{python}{\codedir/classes/NetIndividual.py}

Dopasowanie liczby elementów dzieje
się w~metodzie \texttt{\_update}.
Jest to podzielone na dwa etapy.
Pierwszy to identyfikacja zbyt krótkich
fenotypów i~bezpośrednie skrócenie zbyt długich.
\inputminted[
  firstline=101,
  lastline=113,
  % autogobble,
]{python}{\codedir/classes/NetIndividual.py}

Drugi to wydłużenie zidentyfikowanych fenotypów.
\inputminted[
  firstline=115,
  lastline=122,
  % autogobble,
]{python}{\codedir/classes/NetIndividual.py}

\subsection{%
  \refbykey{sec:mutate}{name}%
}
Decyzja o~mutacji jest podejmowana względem osobnika:
\inputminted[
  firstline=151,
  lastline=154,
  % autogobble,
]{python}{\codedir/classes/Population.py}

Następnie występuje losowanie w~ilu oraz których
podosobnikach nastąpi mutacja:
\inputminted[
  firstline=91,
  lastline=96,
  % autogobble,
]{python}{\codedir/classes/NetIndividual.py}

W~podosobniku następuje kolejne losowanie,
który bit zostanie zmieniony i~on ma zamienioną wartość.
Bit o~wartości $1$ staje się $0$ i~odwrotnie.

Dzieje się to tak samo dla podosobnika sekwencyjnego:
\inputminted[
  firstline=24,
  lastline=26,
  % autogobble,
]{python}{\codedir/classes/GenIndividual.py}

Jak i~listowego:
\inputminted[
  firstline=49,
  lastline=52,
  % autogobble,
]{python}{\codedir/classes/ListIndividual.py}

\subsection{%
  \refbykey{sec:fittness}{name}%
}
W~podrozdziale dotyczącym implementacji klasy pokoleń
zostało pokazane jak realizuje się kilku-krotne uczenie.
Funkcja korygująca w~kodzie ma postać:
\inputminted[
  firstline=17,
  lastline=23,
  % autogobble,
]{python}{\codedir/core.py}

\subsection{%
  \refbykey{sec:fenotype}{name}%
}
Zmiana kodu binarnego na wartość liczbową,
która jest przedstawiona przez
\crefrange{math:first_greater_bin}{math:bin_value},
jest realizowana przez poniższe funkcje:
\inputminted[
  firstline=38,
  lastline=46,
  % autogobble,
]{python}{\codedir/classes/_utils.py}

\subsubsection{%
  \refbykey{subsec:fenotype_seq}{name}%
}
Uzyskanie fenotypu w~podosobniku sekwencyjnym polega
na odczytaniu kolejno części kodu binarnego
respektując długości poszczególnych parametrów:
\inputminted[
  firstline=47,
  lastline=57,
  % autogobble,
]{python}{\codedir/classes/MaxIntsIndividual.py}

\subsubsection{%
  \refbykey{subsec:fenotype_list_nnum}{name}
  oraz \refbykey{subsec:fenotype_list_tactiv}{name}%
}
Uzyskanie fenotypu w~podosobniku listowym polega
na odczytaniu kolejno części kodu binarnego
respektując długość elementów w~danym podosobniku:
\inputminted[
  firstline=47,
  lastline=56,
  % autogobble,
]{python}{\codedir/classes/MaxIntsListIndividual.py}

\subsubsection{Fenotyp -- wyjątki}
Wyjątek stanowią elementy listowych
podosobników oraz parametry sekwencyjnych,
których końcowa oczekiwana
wartość nie jest liczbą całkowitą.
Te które zaliczają się do tej grupy to optymalizator,
wskaźnik uczenia się oraz cała lista typów aktywacji.
Ich wartości na poziomie podfenotypu
zostają w~postacji liczb całkowitych,
dopiero podczas tworzenia sieci neuronowej.

\subsubsection{%
  \refbykey{subsec:fenotype_ind}{name}%
}
Stworzenie sieci neuronowej,
która jest fenotypem osobnika,
dzieje się przez wywołanie
poniższej funkcji podając osobnika,
liczbę parametrów wejściowych i~wyjściowych
oraz rodzaj rozwiązywanej problematyki przez sieć docelową.
\inputminted[
  firstline=77,
  lastline=82,
  % autogobble,
]{python}{\codedir/utils.py}

Funkcja ta zwraca model sieci
oraz parametry uczenia odczytane z~osobnika.
\inputminted[
  firstline=102,
  lastline=106,
  % autogobble,
]{python}{\codedir/utils.py}

By powstał sam model,
sprawdzany jest wpierw typ problematyki,
jeżeli jest to klasyfikacja,
to ostatniej warstwie (wyjściowej)
zostaje przypisany \texttt{softmax}:
\inputminted[
  firstline=57,
  lastline=58,
  % autogobble,
]{python}{\codedir/utils.py}

Następnie tworzony jest szkic struktury modelu,
wraz z~odpowiednimi typami funkcji aktywacyji:
\inputminted[
  firstline=60,
  lastline=63,
  % autogobble,
]{python}{\codedir/utils.py}

A~ostatecznie model jest tworzony z~wskazanym
optymalizatorem oraz wskaźnikiem uczenia się.
\inputminted[
  firstline=64,
  lastline=70,
  % autogobble,
]{python}{\codedir/utils.py}
